\documentclass[a4paper,12pt]{article}
\usepackage[T2A]{fontenc}
\usepackage[utf8]{inputenc}
\usepackage[russian]{babel}
\usepackage{amsmath,amssymb,mathtools}
\usepackage{helvet}\renewcommand{\familydefault}{\sfdefault}
\usepackage{icomma,geometry,microtype,setspace,parskip}
\usepackage{tikz,pgfplots}\usetikzlibrary{positioning}\pgfplotsset{compat=1.18}
\usepackage{tabularx,booktabs,enumitem,xcolor,fancyhdr,hyperref,graphicx}
\geometry{left=2cm,right=2cm,top=2.1cm,bottom=2.1cm}\onehalfspacing
\definecolor{accent}{HTML}{8B3038}\definecolor{teal}{HTML}{287F7A}\definecolor{textgray}{HTML}{30343A}\definecolor{soft}{HTML}{F4EEE5}
\color{textgray}\hypersetup{colorlinks=true,linkcolor=accent,urlcolor=teal}
\setlength{\headheight}{15pt}\pagestyle{fancy}\fancyhf{}\fancyhead[L]{\small Повторение ключевых тем ОГЭ}\fancyhead[R]{\small София Кальней}\fancyfoot[C]{\thepage}
\newcommand{\sectionline}{\par\noindent\textcolor{accent}{\rule{\linewidth}{0.6pt}}\par}
\newcommand{\important}[1]{\textbf{\textcolor{accent}{#1}}}
\begin{document}
\begin{titlepage}\thispagestyle{empty}\centering\vspace*{2.2cm}{\Large\textcolor{accent}{\textbf{ЧЕК-ЛИСТ}}}\par\vspace{.8cm}{\Huge\bfseries Повторение ключевых тем ОГЭ}\par\vspace{.5cm}{\Large Степени, уравнения, неравенства, текстовые задачи и график прямой}\par\vspace{2cm}{\large Лёвин Артём Александрович\par\textcolor{accent}{эксклюзивно для Софии Кальней}}\vfill{\large 29 июля 2026 года}\begin{tikzpicture}[remember picture,overlay]\draw[accent!50,line width=1.2pt]([xshift=1.5cm,yshift=1.5cm]current page.south west)..controls([xshift=6cm,yshift=2.4cm]current page.south west)and([xshift=-6cm,yshift=.7cm]current page.south east)..([xshift=-1.5cm,yshift=1.5cm]current page.south east);\end{tikzpicture}\end{titlepage}
\section*{Главная стратегия}\sectionline
\begin{enumerate}[leftmargin=2em]
\item Определите тип задачи и искомую величину.
\item Выпишите ограничения до преобразований.
\item Выберите модель: формула, таблица, числовая ось или график.
\item Решите задачу школьным методом.
\item Проверьте найденные значения в исходном условии и по смыслу.
\end{enumerate}
\begin{center}\begin{tikzpicture}[node distance=8mm]\node[draw=accent,rounded corners,inner sep=7pt](a){Распознать тип};\node[draw=teal,rounded corners,inner sep=7pt,right=of a](b){Записать ограничения};\node[draw=accent,rounded corners,inner sep=7pt,right=of b](c){Выбрать модель};\node[draw=teal,rounded corners,inner sep=7pt,below=of b](d){Решить и проверить};\draw[->,thick](a)--(b);\draw[->,thick](b)--(c);\draw[->,thick](b)--(d);\end{tikzpicture}\end{center}
\clearpage
\section*{1. Степени}\sectionline
\[ (a^m)^n=a^{mn},\qquad (ab)^n=a^nb^n,\qquad a^ma^n=a^{m+n},\qquad \frac{a^m}{a^n}=a^{m-n},\ a\ne0. \]
\[ a^{-n}=\frac1{a^n}. \]
\important{Алгоритм:} раскрыть степень степени; объединить одинаковые основания; сократить; сохранить ограничения.
\[\frac{(x^3)^2y^4}{x^{11}(y^3)^2}=x^{6-11}y^{4-6}=\frac1{x^5y^2},\qquad x\ne0,\ y\ne0.\]
\section*{2. Уравнение с квадратным корнем}\sectionline
\[x=\sqrt{x+6}.\]
Поскольку правая часть неотрицательна, \(x\ge0\). После возведения в квадрат:
\[x^2=x+6,\qquad (x-3)(x+2)=0.\]
Кандидаты: \(3\) и \(-2\). Значение \(-2\) не удовлетворяет условию \(x\ge0\). Проверка \(3=\sqrt9\) верна, поэтому \(x=3\).
\clearpage
\section*{3. Система линейных неравенств}\sectionline
\[\begin{cases}-8+4x>0,\\4-3x>-8.\end{cases}\]
Получаем \(x>2\) и \(x<4\). Для системы берётся пересечение:
\[x\in(2;4).\]
\begin{center}\begin{tikzpicture}[x=1.5cm]\draw[->](-1,0)--(5,0);\draw[accent,line width=2.5pt](2.08,0)--(3.92,0);\draw[accent,fill=white,line width=1pt](2,0)circle(2.5pt)(4,0)circle(2.5pt);\node[below]at(2,0){2};\node[below]at(4,0){4};\end{tikzpicture}\end{center}
\important{Помните:} при делении неравенства на отрицательное число знак меняется.
\section*{4. Метод интервалов}\sectionline
\[\frac{-12}{x^2-7x-8}<0,\qquad x^2-7x-8=(x+1)(x-8).\]
Критические точки \(-1\) и \(8\) являются нулями знаменателя и всегда выколоты. Знаки дроби: минус, плюс, минус. Ответ:
\[x\in(-\infty;-1)\cup(8;+\infty).\]
\begin{center}\begin{tikzpicture}[x=1.25cm]\draw[->](-4,0)--(4.5,0);\draw[accent,fill=white,line width=1pt](-2,0)circle(2.5pt)(2,0)circle(2.5pt);\node[below]at(-2,0){$-1$};\node[below]at(2,0){$8$};\node[above]at(-3,0){$-$};\node[above]at(0,0){$+$};\node[above]at(3,0){$-$};\draw[teal,line width=2.5pt](-3.8,0)--(-2.08,0)(2.08,0)--(4.1,0);\end{tikzpicture}\end{center}
\clearpage
\section*{5. Растворы}\sectionline
Масса растворённого вещества:
\[m_{\text{вещества}}=m_{\text{раствора}}\cdot c.\]
Пусть \(x\) и \(y\) — концентрации двух растворов массами 10 кг и 16 кг. Из смешивания всего содержимого:
\[10x+16y=26\cdot0{,}55.\]
Из смешивания равных масс:
\[x+y=1{,}22.\]
Решение системы даёт \(y=0{,}35\), \(x=0{,}87\). Следовательно, масса кислоты в первом растворе:
\[10x=8{,}7\text{ кг}.\]
\section*{6. Движение и работа}\sectionline
Для движения \(S=vt\), \(t=S/v\). Катер, прошедший по 35 км по и против течения со скоростью течения 4 км/ч за 3 часа, удовлетворяет уравнению
\[\frac{35}{x+4}+\frac{35}{x-4}=3,\qquad x>4.\]
После решения получаем \(x=24\) км/ч.

Для работы \(A=pt\), \(t=A/p\). Если первый рабочий производит на 10 деталей в час больше и выполняет заказ из 60 деталей на 3 часа быстрее, то
\[\frac{60}{x}-\frac{60}{x+10}=3,\qquad x>0.\]
Подходит корень \(x=10\) деталей в час.
\clearpage
\section*{7. График прямой и параметр}\sectionline
Линейная функция имеет вид \(y=kx+b\). Коэффициент \(k\) задаёт наклон, \(b\) — пересечение с осью \(Oy\).

В семействе
\[y=x+a-3\]
наклон постоянен и равен 1, поэтому прямые параллельны и сдвигаются вверх или вниз.

В семействе
\[y=(a-3)x+1\]
свободный коэффициент постоянен, поэтому все прямые проходят через точку \((0;1)\) и меняют наклон.
\begin{center}\begin{tikzpicture}\begin{axis}[width=.8\linewidth,height=7cm,axis lines=middle,unit vector ratio=1 1,xmin=-4,xmax=4,ymin=-4,ymax=4,samples=220,grid=both,clip=false]\addplot[accent,thick,domain=-4:4]{x-2};\addplot[teal,thick,domain=-4:4]{x};\addplot[accent!60,thick,domain=-4:4]{x+2};\end{axis}\end{tikzpicture}\end{center}
\clearpage
\section*{Тренировочный блок}\sectionline
\begin{enumerate}[leftmargin=2em,itemsep=.8em]
\item Упростите \(\dfrac{(a^5b^2)^3}{a^{11}b^4}\).
\item Решите \(x=\sqrt{x+12}\).
\item Решите систему \(x>-1\), \(x\le5\).
\item Решите \(\dfrac{x+5}{x-3}\ge0\).
\item Решите \(\dfrac{-18}{x^2-5x-14}<0\).
\item Смешали 8 кг 20\%-го и 12 кг 50\%-го растворов. Найдите концентрацию смеси.
\item Катер прошёл 35 км по и против течения за 3 часа при скорости течения 4 км/ч. Найдите собственную скорость.
\item Первый рабочий производит на 10 деталей в час больше и выполняет заказ из 60 деталей на 3 часа быстрее. Найдите производительность второго.
\item Опишите семейство прямых \(y=x+a-3\).
\item При каком \(a\) прямые \(y=x+a-3\) и \(y=(a-3)x+1\) перпендикулярны?
\end{enumerate}
\clearpage
\section*{Ответы}\sectionline
\renewcommand{\arraystretch}{1.35}\begin{tabularx}{\linewidth}{cX}\toprule №&Ответ\\\midrule1&\(a^4b^2\)\\2&\(x=4\)\\3&\((-1;5]\)\\4&\((-\infty;-5]\cup(3;+\infty)\)\\5&\((-\infty;-2)\cup(7;+\infty)\)\\6&38\%\\7&24 км/ч\\8&10 деталей/ч\\9&Параллельные прямые с наклоном 1\\10&\(a=2\)\\\bottomrule\end{tabularx}
\vfill\begin{center}\textcolor{accent}{\textbf{Финальная привычка: ограничения → решение → подстановка → ответ по вопросу.}}\end{center}
\end{document}
